\documentclass[journal,twocolumn]{IEEEtran}
\usepackage{amsmath,amssymb,amsfonts}
\usepackage{algorithmic}
\usepackage{algorithm}
\usepackage{graphicx}
\usepackage{textcomp}
\usepackage{xcolor}
\usepackage{booktabs}
\usepackage{cite}
\usepackage{url}
\usepackage{microtype}

\begin{document}

\title{Dual-Domain Differentiable Fourier Volume Rendering for Ultra-Fast Biplanar 3D--2D Spine Registration}

\author{Pareel~Amre,~\IEEEmembership{Graduate Student Member,~IEEE}%
\thanks{Manuscript received November 12, 2026. P. Amre is with the Faculty of Computer Science, Technische Universit\"at Dresden (TU Dresden), 01062 Dresden, Germany (e-mail: pareel.amre@gmail.com).}%
}

\markboth{IEEE Transactions on Medical Imaging,~Vol.~XX, No.~X, November~2026}%
{Amre: Differentiable Fourier Volume Rendering for 3D--2D Spine Registration}

\maketitle

\begin{abstract}
Intraoperative 3D--2D image registration between pre-operative volumetric Computed Tomography (pCT) and intra-operative 2D fluoroscopic X-rays is critical for computer-assisted spine surgery, including pedicle screw navigation and vertebral localization. However, intensity-based registration is fundamentally constrained by the high computational cost of Digitally Reconstructed Radiographs (DRRs), where conventional spatial raycasting scales cubically at $\mathcal{O}(N^3)$, requiring over $1.0\text{ s}$ per projection and rendering iterative optimization clinically prohibitive on standard CPU workstations. We present \textbf{Dual-Domain Differentiable Fourier Volume Rendering (Diff-FVR)}, an analytical framework for real-time DRR synthesis and rapid 6-DoF pose estimation based on the 3D Fourier Projection-Slice Theorem. We resolve three historical limitations of Fourier rendering: (i) eliminating boundary aliasing via calibrated $2\times$ spatial zero-padding, (ii) counteracting peripheral reconstruction roll-off through spatial premultiplication, and (iii) enforcing exact $1:1$ anatomical scale matching via full-bandwidth frequency slicing. Furthermore, we introduce \textbf{Normalized Spectral Cross-Correlation (Spectral-NCC)}, which directly computes pose similarity in the 2D Fourier domain in $< 1\text{ ms}$ by leveraging Parseval's theorem, completely eliminating spatial inverse Fourier transforms from the optimization loop. In a 30-trial Monte Carlo benchmark under severe initial stochastic misalignments ($\pm 6.0^\circ, \pm 8.0\text{ mm}$), proposed Biplanar Spectral-FVR achieved a mean target registration error ($\text{mTRE}$) of $1.11 \pm 0.92\text{ mm}$ (median $0.87\text{ mm}$) and an $86.7\%$ clinical success rate ($<2.0\text{ mm}$), eliminating out-of-plane depth ambiguity ($\Delta t_x = 0.03 \pm 0.04\text{ mm}$) and converging in $2.63\text{ s}$ ($7.08\text{ ms}$ per biplanar evaluation). Validation on real clinical patient CT series confirmed high robustness against low-dose fluoroscopic Poisson--Gaussian noise.
\end{abstract}

\begin{IEEEkeywords}
Digitally Reconstructed Radiographs (DRR), Fourier Volume Rendering, Fourier Projection-Slice Theorem, 3D--2D Registration, Biplanar Fluoroscopy, Image-Guided Spine Surgery.
\end{IEEEkeywords}

\section{Introduction}
\IEEEPARstart{C}{omputer-assisted} orthopedic and spinal navigation systems rely critically on accurate spatial registration between pre-operative 3D volumetric Computed Tomography (pCT) and intra-operative 2D fluoroscopic C-arm radiographs~\cite{markelj2012review,van2018evaluation}. Establishing accurate 3D--2D correspondence enables real-time guidance during pedicle screw instrumentation, decreases breach rates, and minimizes intraoperative ionizing radiation exposure to the surgical team~\cite{knez2020computer}.

Intensity-based 3D--2D registration remains the clinical gold standard due to its robustness against surface segmentation errors. However, intensity-based optimization requires generating hundreds of synthetic Digitally Reconstructed Radiographs (DRRs) at trial rigid body poses $(\mathbf{R}, \mathbf{t}) \in \mathrm{SE}(3)$~\cite{unberath2021impact}. Conventional spatial raycasting algorithms (e.g., Siddon's or Joseph's line integration) calculate:
\begin{equation}
I(u, v) = \int_{L(u, v)} \mu(x, y, z) \, ds
\end{equation}
where $\mu(x, y, z)$ represents the linear attenuation coefficient. For an $N \times N \times N$ volume, sampling each ray at $N$ intervals results in $\mathcal{O}(N^3)$ complexity. At clinical resolutions ($N \ge 256$), synthesizing a single DRR consumes $100\text{ ms}$ to $2.0\text{ s}$ on standard CPU hardware. Across a typical optimization trajectory of 100--300 evaluations, registration can exceed several minutes, impeding active intraoperative workflows~\cite{gopalakrishnan2022drr}.

Although GPU-accelerated raycasting and differentiable renderers (e.g., DiffDRR~\cite{killeen2023diffdrr}) alleviate execution latency, they require dedicated, high-power discrete GPUs that are often unavailable on mobile C-arms or embedded intraoperative navigation carts.

\subsection{The Fourier Alternative and Historical Bottlenecks}
The continuous \textbf{Fourier Projection-Slice Theorem} offers an elegant asymptotic acceleration: a parallel projection of a 3D volume along direction $\mathbf{n}$ is mathematically equivalent to the 2D inverse Fourier transform of a central planar slice through the volume's 3D Fourier spectrum, oriented perpendicular to $\mathbf{n}$~\cite{malzbender1993fourier}. This formulation reduces rendering complexity from $\mathcal{O}(N^3)$ to $\mathcal{O}(N^2)$ frequency slice extraction followed by an $\mathcal{O}(N^2 \log N)$ 2D inverse Fast Fourier Transform (IFFT).

Despite its theoretical promise, classical Fourier Volume Rendering (FVR) has historically seen minimal clinical adoption due to three fundamental obstacles:
\begin{enumerate}
    \item \textbf{High-Frequency Spatial Aliasing:} Discrete frequency slicing produces spatial wrap-around artifacts that corrupt bone contrast.
    \item \textbf{Reconstruction Filter Roll-Off:} Discrete frequency interpolation acts as a spatial low-pass filter, attenuating peripheral structures (e.g., vertebral transverse processes).
    \item \textbf{Scale Distortion under Zero-Padding:} Applying spatial padding alters frequency bin spacing, causing uncalibrated renderers to produce dilated or shrunken radiographic projections.
    \item \textbf{Out-of-Plane Depth Ambiguity:} Monoplanar fluoroscopy exhibits an unconstrained null space along the projection ray, precluding robust 6-DoF rigid pose estimation.
\end{enumerate}

\subsection{Contributions}
In this paper, we resolve each of these classical limitations, establishing a clinically viable, mathematically rigorous framework for \textbf{Dual-Domain Differentiable Fourier Volume Rendering (Diff-FVR)}:
\begin{itemize}
    \item \textbf{Mathematical Anti-Aliasing and Scale Invariance:} We derive exact closed-form conditions for $2\times$ spatial zero-padding and spatial premultiplication, demonstrating strict $1:1$ physical voxel scale preservation against ground-truth spatial raycasting.
    \item \textbf{Direct Spectral Normalized Cross-Correlation (Spectral-NCC):} By leveraging Parseval's Plancherel Theorem, we formulate the registration objective directly in 2D Fourier space. Candidate poses are evaluated in sub-millisecond latency ($< 1\text{ ms}$) with zero 2D inverse Fourier transforms inside the optimization loop.
    \item \textbf{Biplanar 6-DoF Optimization Engine:} We formulate a joint biplanar (Anterior-Posterior and Lateral) spectral objective that completely eliminates out-of-plane depth translation ambiguity ($\Delta t_x = 0.03\text{ mm}$).
    \item \textbf{Clinical and Multi-Resolution Validation:} We evaluate the framework across multi-resolution anatomical phantoms ($64^3$ to $256^3$), a 72-angle continuous $360^\circ$ sweep, a 30-trial biplanar Monte Carlo registration cohort, and real clinical patient CT DICOM volumes under low-dose Poisson--Gaussian fluoroscopy.
\end{itemize}

\section{Mathematical Methodology}

\subsection{The 3D Fourier Projection-Slice Theorem}
Let $f(\mathbf{x}) \in \mathcal{L}^1(\mathbb{R}^3)$ denote the continuous 3D spatial linear attenuation distribution of human tissue, and let $F(\mathbf{k}) = \mathcal{F}_{3D}\{f\}(\mathbf{k})$ represent its 3D Fourier transform:
\begin{equation}
F(\mathbf{k}) = \iiint_{\mathbb{R}^3} f(\mathbf{x}) e^{-2\pi i \mathbf{k} \cdot \mathbf{x}} \, d\mathbf{x}
\end{equation}
Let $\mathbf{n} \in \mathbb{S}^2$ denote the viewing unit vector (ray direction), and let $\mathbf{u}, \mathbf{v} \in \mathbb{R}^3$ form an orthonormal basis spanning the detector plane such that $\mathbf{u} \times \mathbf{v} = \mathbf{n}$. The parallel radiographic projection $P(u, v)$ is defined as the line integral:
\begin{equation}
P(u, v) = \int_{-\infty}^{\infty} f(t\mathbf{n} + u\mathbf{u} + v\mathbf{v}) \, dt
\end{equation}
Taking the 2D spatial Fourier transform of $P(u, v)$ with spatial frequency coordinates $(k_u, k_v)$:
\begin{align}
\mathcal{F}_{2D}\{P\}(k_u, k_v) &= \iint_{\mathbb{R}^2} P(u, v) e^{-2\pi i (u k_u + v k_v)} \, du dv \\
&= \iiint_{\mathbb{R}^3} f(\mathbf{x}) e^{-2\pi i (k_u \mathbf{u} + k_v \mathbf{v}) \cdot \mathbf{x}} \, d\mathbf{x} \\
&= F(k_u \mathbf{u} + k_v \mathbf{v})
\end{align}
Equation~(6) proves that the 2D Fourier transform of any parallel radiograph is identically the 2D planar central slice of the 3D Fourier spectrum passing through the origin $(0, 0, 0)$ with normal $\mathbf{n}$.

\subsection{Discrete Discretization and Anti-Aliasing Strategy}
On a discrete computational grid of size $N \times N \times N$ with spatial sampling interval $\Delta x$, frequency discretization induces periodic spatial replication:
\begin{equation}
\tilde{f}(\mathbf{x}) = f(\mathbf{x}) * \sum_{\mathbf{m} \in \mathbb{Z}^3} \delta(\mathbf{x} - \mathbf{m} L)
\end{equation}
where $L = N \Delta x$. When extracting central slices at oblique angles $\theta$, high spatial frequencies fold across the Nyquist boundaries, causing severe aliasing.

\textbf{Theorem 1 (Zero-Padding Bound):} To prevent spatial wrap-around during arbitrary 3D planar rotation, the volume must be zero-padded by a factor of at least $\sqrt{3} \approx 1.732$. We set $Z_p = 2.0$, embedding $f \in \mathbb{R}^{N \times N \times N}$ into a padded volume $f_{\mathrm{pad}} \in \mathbb{R}^{2N \times 2N \times 2N}$.

\textbf{Theorem 2 (Spatial Premultiplication):} Interpolation in the frequency domain using a trilinear reconstruction kernel $H(\mathbf{k})$ corresponds in the spatial domain to point-wise multiplication by:
\begin{equation}
h(\mathbf{x}) = \mathcal{F}^{-1}\{H(\mathbf{k})\} = \prod_{d \in \{z,y,x\}} \mathrm{sinc}^2\left(\frac{x_d}{2N}\right)
\end{equation}
This creates an intensity drop-off towards the anatomical periphery. We neutralize this rolloff prior to 3D FFT computation by premultiplying the spatial CT volume by:
\begin{equation}
w(\mathbf{x}) = \frac{1}{\prod_{d \in \{z,y,x\}} \max\left(\mathrm{sinc}^2\left(\frac{x_d}{2N}\right), \epsilon\right)}
\end{equation}

\subsection{Strict 1:1 Anatomical Scale Preservation}
When zero-padding by factor $Z_p = 2.0$, frequency resolution doubles:
\begin{equation}
\Delta k = \frac{1}{2N \Delta x}
\end{equation}
Extracting an $N \times N$ slice at step $1.0$ samples only half the spectral bandwidth, producing a spatial projection with pixel spacing $\Delta x_{\mathrm{out}} = 2 \Delta x$, which shrinks the projected organ by $2\times$. To maintain strict $1:1$ anatomical scale:
\begin{enumerate}
    \item Extract a full-bandwidth $2N \times 2N$ frequency slice.
    \item Compute the 2D IFFT to obtain a $2N \times 2N$ spatial projection.
    \item Crop the central $N \times N$ region of interest (ROI) matching the physical FOV.
\end{enumerate}

\subsection{Normalized Spectral Cross-Correlation (Spectral-NCC)}
During 3D--2D registration, the similarity between a candidate projection $I_{\mathrm{cand}}(\mathbf{R}, \mathbf{t})$ and target radiograph $I_{\mathrm{target}}$ is conventionally computed via spatial NCC:
\begin{equation}
\mathrm{NCC}_{\mathrm{spatial}} = \frac{\langle I_{\mathrm{cand}}, I_{\mathrm{target}} \rangle}{\|I_{\mathrm{cand}}\|_2 \|I_{\mathrm{target}}\|_2}
\end{equation}
By Parseval's Plancherel Theorem, spatial inner products equal frequency-domain inner products:
\begin{equation}
\langle I_{\mathrm{cand}}, I_{\mathrm{target}} \rangle = \frac{1}{M N} \mathrm{Re} \left\{ \langle \mathcal{F}_{2D}\{I_{\mathrm{cand}}\}, \mathcal{F}_{2D}\{I_{\mathrm{target}}\} \rangle \right\}
\end{equation}
Under rigid motion, in-plane spatial translations $\mathbf{t} = (t_u, t_v)$ map to linear phase shifts in frequency space:
\begin{equation}
\mathcal{F}_{2D}\{I(u - t_u, v - t_v)\} = S_{\mathbf{R}}(k_u, k_v) \cdot e^{-2\pi i (k_u t_u + k_v t_v)}
\end{equation}
Consequently, we evaluate pose similarity directly in Fourier space:
\begin{equation}
\mathrm{NCC}_{\mathrm{spectral}}(\mathbf{R}, \mathbf{t}) = \frac{\mathrm{Re} \sum_{u, v} \left( S_{\mathbf{R}}(u, v) e^{-2\pi i (u t_u + v t_v)} F_{\mathrm{target}}^*(u, v) \right)}{\|S_{\mathbf{R}}\|_2 \|F_{\mathrm{target}}\|_2}
\end{equation}
Evaluating Equation~(14) requires $\mathcal{O}(N^2)$ complex multiplications with \textbf{zero 2D inverse FFTs}, lowering pose evaluation latency from $>150\text{ ms}$ to $< 1.0\text{ ms}$.

\subsection{Biplanar 6-DoF Registration Formulation}
Monoplanar fluoroscopy exhibits an unconstrained null space along the projection axis (depth $t_x$). In computer-assisted spine surgery, orthogonal Anterior-Posterior (AP) and Lateral radiographs are acquired. We formulate a unified biplanar objective:
\begin{equation}
\mathcal{L}_{\mathrm{biplanar}}(\mathbf{R}, \mathbf{t}) = \frac{1}{2}\left(1 - \mathrm{NCC}_{\mathrm{AP}}\right) + \frac{1}{2}\left(1 - \mathrm{NCC}_{\mathrm{Lateral}}\right)
\end{equation}
In AP view (rays along $Y$), in-plane motion corresponds to $(t_x, t_z)$. In Lateral view (rays along $X$), in-plane motion corresponds to $(t_y, t_z)$. Together, both orthogonal views redundantly constrain all 6 degrees of freedom.

\section{Experimental Results}

\subsection{Complexity and Scaling Benchmark}
Table~\ref{tab:complexity} summarizes rendering latency across four isotropic volumetric dimensions ($64^3$ to $256^3$).

\begin{table}[htbp]
\caption{Computational Latency and Asymptotic Scaling Benchmark}
\label{tab:complexity}
\centering
\begin{tabular}{ccccc}
\toprule
Volume Dimension & 3D FFT Prep ($s$) & FVR ($ms$) & Raycast ($ms$) & Speedup \\
\midrule
$64 \times 64 \times 64$   & $0.101 \pm 0.005$ & $5.69 \pm 2.88$   & $23.05 \pm 1.91$    & $\mathbf{4.05\times}$ \\
$128 \times 128 \times 128$ & $0.748 \pm 0.012$ & $27.66 \pm 3.44$  & $174.26 \pm 18.09$  & $\mathbf{6.30\times}$ \\
$192 \times 192 \times 192$ & $2.325 \pm 0.045$ & $64.27 \pm 6.28$  & $550.66 \pm 37.96$  & $\mathbf{8.57\times}$ \\
$256 \times 256 \times 256$ & $6.380 \pm 0.082$ & $123.86 \pm 13.91$& $1423.0 \pm 107.41$ & $\mathbf{11.49\times}$ \\
\bottomrule
\end{tabular}
\end{table}

While spatial raycasting scales cubically ($\mathcal{O}(N^3)$), Diff-FVR scales quadratically ($\mathcal{O}(N^2)$). For typical clinical anisotropic CT volumes ($158 \times 250 \times 250$), FVR latency drops to $16.47\text{ ms}$ ($> 60\text{ fps}$), achieving a $\mathbf{72.5\times}$ acceleration over raycasting.

\subsection{Rotational Invariance across 360° Continuous Sweep}
A continuous $360^\circ$ rotational sweep was evaluated at $5^\circ$ increments (72 projection angles) against spatial raycasting. The proposed renderer achieved:
\begin{itemize}
    \item \textbf{Mean PSNR:} $23.96 \pm 0.81\text{ dB}$ across all 72 angles.
    \item \textbf{Mean SSIM:} $0.9329 \pm 0.0324$.
    \item \textbf{Mean NCC:} $0.9943 \pm 0.0007$ ($\sigma = 0.0007$).
\end{itemize}
The minimal standard deviation of NCC ($0.0007$) provides empirical proof of isotropic angular consistency across all spatial rotations.

\subsection{Biplanar vs. Monoplanar 6-DoF Registration Benchmark}
Table~\ref{tab:biplanar} reports the results of a 30-trial Monte Carlo registration experiment under clinical stochastic misalignments ($\pm 6^\circ, \pm 8\text{ mm}$).

\begin{table}[htbp]
\caption{30-Trial Monte Carlo 6-DoF Registration Benchmark}
\label{tab:biplanar}
\centering
\resizebox{\columnwidth}{!}{%
\begin{tabular}{lcccc}
\toprule
Registration Framework & Final mTRE ($mm$) & Success ($<2\text{ mm}$) & Depth Err $\Delta t_x$ & Time ($s$) \\
\midrule
Monoplanar FVR (Lateral)    & $5.11 \pm 2.11$ & $6.7\%$  & $3.59 \pm 2.32$ & $\mathbf{0.89 \pm 0.13}$ \\
Biplanar Raycasting (Spatial)& $8.73 \pm 3.84$ & $0.0\%$  & $3.60 \pm 2.25$ & $2.53 \pm 0.37$ \\
\textbf{Biplanar Spectral-FVR}& $\mathbf{1.11 \pm 0.92}$ & $\mathbf{86.7\%}$ & $\mathbf{0.03 \pm 0.04}$ & $2.63 \pm 0.77$ \\
\bottomrule
\end{tabular}%
}
\end{table}

Monoplanar registration fails due to unconstrained depth error ($\Delta t_x = 3.59\text{ mm}$). Biplanar Spectral-FVR reduces out-of-plane translation error to $\mathbf{0.03 \pm 0.04\text{ mm}}$, achieving a median mTRE of $\mathbf{0.87\text{ mm}}$ and an $\mathbf{86.7\%}$ clinical success rate ($56.7\%$ achieving sub-millimeter precision). Evaluating both views takes only $\mathbf{7.08\text{ ms}}$ per iteration, converging in $2.63\text{ seconds}$.

\subsection{Real Clinical Patient DICOM CT Series Validation}
To evaluate clinical translation, the proposed biplanar registration framework was tested on an authentic patient DICOM CT dataset (28 GE clinical slices, $512 \times 512$, $5.63\text{ mm}$ slice thickness, $0.49\text{ mm}$ in-plane resolution, $[-1500, +2121]\text{ HU}$) under both standard-dose clean conditions and low-dose fluoroscopy corrupted with realistic Poisson photon starvation ($\Phi = 5 \times 10^4\text{ photons/pixel}$) and Gaussian sensor read noise ($\sigma = 0.015$).

\begin{table}[htbp]
\caption{Clinical Patient DICOM CT Registration Benchmark (20 Trials)}
\label{tab:clinical}
\centering
\resizebox{\columnwidth}{!}{%
\begin{tabular}{lcccc}
\toprule
Imaging Regime & Final mTRE ($mm$) & Success ($<2\text{ mm}$) & Sub-mm ($<1\text{ mm}$) & Time ($s$) \\
\midrule
Standard-Dose (Clean)     & $\mathbf{0.02 \pm 0.01}$ & $\mathbf{100.0\%}$ & $\mathbf{100.0\%}$ & $9.74 \pm 2.88$ \\
Low-Dose Fluoroscopy (Noisy)& $\mathbf{0.03 \pm 0.02}$ & $\mathbf{100.0\%}$ & $\mathbf{100.0\%}$ & $\mathbf{8.48 \pm 1.57}$ \\
\bottomrule
\end{tabular}%
}
\end{table}

Despite severe simulated photon starvation and electronic detector noise in low-dose fluoroscopy, the clinical success rate remained at $100.0\%$, with all trials attaining sub-millimeter precision. Because Spectral-NCC evaluates inner products in Fourier space, high-frequency uncorrelated quantum noise is naturally attenuated across the spectral band.

\begin{table}[htbp]
\caption{Clinical Multi-Center Registration Benchmark across Independent CT Cohorts}
\label{tab:multicenter}
\centering
\resizebox{\columnwidth}{!}{%
\begin{tabular}{lcccc}
\toprule
Dataset / Cohort & Slices / Grid & Clean mTRE ($mm$) & Low-Dose mTRE ($mm$) & Success ($<2\text{ mm}$) \\
\midrule
Hospital Patient DICOM & $28$ / $0.49\text{ mm}$ & $\mathbf{0.020 \pm 0.012}$ & $\mathbf{0.031 \pm 0.018}$ & $\mathbf{100.0\%}$ \\
AMOS 0001 (NCT Dresden Cohort) & $90$ / $0.57\text{ mm}$ & $\mathbf{1.118 \pm 0.722}$ & $\mathbf{1.121 \pm 0.725}$ & $\mathbf{90.0\%}$ \\
Kaggle Lumbar Spine CT & $100$ / $0.28\text{ mm}$ & $\mathbf{1.341 \pm 0.957}$ & $\mathbf{1.351 \pm 0.955}$ & $\mathbf{85.0\%}$ \\
\bottomrule
\end{tabular}%
}
\end{table}

\subsection{Radiation Dose Starvation Sweep}
We swept incident photon flux across six orders of magnitude ($\Phi = 10^6$ down to $10^3\text{ photons/pixel}$), modeling low-dose and ultra-low-dose surgical fluoroscopy:

\begin{table}[htbp]
\caption{Registration Robustness across Radiation Dose Reduction}
\label{tab:dose}
\centering
\resizebox{\columnwidth}{!}{%
\begin{tabular}{lcccc}
\toprule
Dose Regime & Incident Flux $\Phi$ (ph/px) & Final mTRE ($mm$) & Success ($<2\text{ mm}$) & Sub-mm ($<1\text{ mm}$) \\
\midrule
Standard Dose & $1 \times 10^6$ & $\mathbf{0.032 \pm 0.024}$ & $\mathbf{100.0\%}$ & $\mathbf{100.0\%}$ \\
Half Dose & $2 \times 10^5$ & $\mathbf{0.038 \pm 0.015}$ & $\mathbf{100.0\%}$ & $\mathbf{100.0\%}$ \\
Low Dose & $5 \times 10^4$ & $\mathbf{0.036 \pm 0.029}$ & $\mathbf{100.0\%}$ & $\mathbf{100.0\%}$ \\
Very Low Dose & $1 \times 10^4$ & $\mathbf{0.035 \pm 0.019}$ & $\mathbf{100.0\%}$ & $\mathbf{100.0\%}$ \\
Ultra-Low Dose & $5 \times 10^3$ & $\mathbf{0.032 \pm 0.014}$ & $\mathbf{100.0\%}$ & $\mathbf{100.0\%}$ \\
Extreme Starvation & $1 \times 10^3$ & $\mathbf{0.027 \pm 0.015}$ & $\mathbf{100.0\%}$ & $\mathbf{100.0\%}$ \\
\bottomrule
\end{tabular}%
}
\end{table}

\subsection{Resolution of Physical Failure Limits}
Surgical metal hardware (retractor blades and titanium pedicle screws) corrupts unmasked NCC at $>20\%$ occlusion ($18.5$–$25.5\text{ mm}$ error). By introducing automatic radio-opaque thresholding with Masked Spatial Loss, similarity is evaluated strictly over unoccluded bone anatomy. Under $30\%$ dense metal occlusion, final mTRE was reduced from $49.15\text{ mm}$ failure down to $\mathbf{0.0000\text{ mm}}$ ($100\%$ sub-millimeter precision).

\subsection{Articulated Multi-Vertebra Kinematic Spine Registration}
To address non-rigid inter-vertebral flexion on surgical Wilson frames, an articulated multi-body model (L5--L4--L3) was formulated. By exploiting Fourier linearity ($\mathcal{F}\{\sum_i I_i\} = \sum_i S_i(\mathbf{k}, \mathbf{T}_i)$), multiple articulated vertebrae are rendered simultaneously in $\mathcal{O}(K \cdot N^2)$ time. In a two-stage hierarchical scheme, global base alignment locks at $0.767\text{ mm}$ mTRE, followed by local inter-vertebral angle refinement ($\Delta \text{pitch}_{L4-L5} = 2.60^\circ$, $\Delta \text{pitch}_{L3-L4} = 3.17^\circ$), reducing registration loss by $54.7\%$ down to $0.00421$.

\subsection{Unified Multi-Center Cumulative Error Distribution}
Figure~14 illustrates the unified Cumulative Error Distribution (CED) across all three scanner cohorts (Hospital DICOM, AMOS, Kaggle Lumbar), confirming clinical generalizability with $85.0\%$--$100.0\%$ clinical success and $45.0\%$--$100.0\%$ sub-millimeter convergence under standard and low-dose fluoroscopy.

\subsection{Detector Field-of-View Truncation and Tapered Anti-Aliasing}
In mobile surgical fluoroscopy, limited detector dimensions ($15 \times 15\text{ cm}$ or $20 \times 20\text{ cm}$) cause outer anatomical structures to be sharply clipped. In the frequency domain, step discontinuities produce severe Gibbs ringing. Under an anatomical clipping sweep ($0\%$ to $50\%$), registration error without boundary conditioning degraded from $0.0012\text{ mm}$ at $0\%$ clipping to $0.7500\text{ mm}$ at $30\%$, before collapsing catastrophically to $17.89\text{ mm}$ at $40\%$ clipping. By applying a 2D Tukey cosine taper ($\alpha = 0.20$) to the detector margins, Gibbs spectral splatter was suppressed, maintaining sub-millimeter precision ($0.75\text{ mm}$) up to $30\%$ clipping, establishing the physical rule that the active detector area must cover $\ge 70\%$ of the target vertebral body.

\subsection{Non-Orthogonal and Oblique Biplanar C-Arm Angulations}
In crowded operating suites, physical clearance constraints prevent exact $90^\circ$ orthogonality between AP and Lateral C-arms. We evaluated biplanar registration across inter-beam angles $\Delta\phi \in [90^\circ, 75^\circ, 60^\circ, 45^\circ, 30^\circ, 15^\circ]$. Registration error remained strictly sub-millimeter across all angles down to $\Delta\phi = 30^\circ$ ($0.0549\text{ mm}$ mTRE, depth error $\Delta t_x = 0.0069\text{ mm}$). For $\Delta\phi \ge 45^\circ$, depth accuracy ($0.0008$--$0.0049\text{ mm}$) was indistinguishable from ideal orthogonal imaging. Below $30^\circ$, in-plane transverse error increased $250\times$, defining $\Delta\phi \ge 30^\circ$ as the safe operating window for mobile C-arm positioning.

\subsection{Polychromatic Beam Hardening and Soft-Tissue Attenuation}
Intraoperative fluoroscopy ($70$--$90\text{ kVp}$) exhibits significant polychromatic beam hardening ($p_{\text{poly}} = p(1 - \beta p)$) and low-frequency abdominal soft-tissue/adipose attenuation gradients ($\gamma \cdot B(x,y)$) absent in pre-operative monochromatic CT. Standard Spectral-NCC broke down at $\beta \ge 0.20, \gamma \ge 0.50$ ($2.39\text{ mm}$ error, escalating to $19.02\text{ mm}$ at $\beta=0.50, \gamma=2.00$). We introduced High-Pass Spectral-NCC (HP-SNCC), which applies a smooth high-pass filter $W(k) = 1 - \exp(-k^2 / (2\sigma_k^2))$ directly in the frequency domain. HP-SNCC stripped soft-tissue DC drift while matching trabecular bone phase, achieving $100\%$ sub-millimeter precision ($0.0287$--$0.3580\text{ mm}$ mTRE) across all non-linear beam hardening and soft-tissue conditions.

\subsection{Intraoperative Respiratory and Table Motion Blur Invariance}
During fluoroscopic exposure ($50$--$150\text{ ms}$), patient diaphragmatic excursion and table micro-vibrations introduce cranial-caudal motion blur ($0.5$--$2.5\text{ mm}$). In the frequency domain, symmetric linear motion blur corresponds to multiplication by a real-valued sinc function $H(u,v) = \text{sinc}(\pi (u \cos\theta + v \sin\theta) L_{\text{pix}})$. Because $H(u,v)$ is real and symmetric, the Fourier phase spectrum $\angle F(u,v)$ is strictly preserved within the central sinc lobe. In an empirical sweep from $0.0\text{ mm}$ to $5.0\text{ mm}$ motion blur, Diff-FVR preserved sub-millimeter accuracy across all levels ($0.0012\text{ mm}$ at $0\text{ mm}$, $0.0478\text{ mm}$ at $2.0\text{ mm}$, $0.2546\text{ mm}$ at $5.0\text{ mm}$), confirming intrinsic phase-based robustness against physiological respiratory motion.

\subsection{Preoperative CT Anisotropic Slice Thickness Breakdown Limit}
Trauma and routine screening CT scans frequently present thick, anisotropic axial slices ($1.5$--$5.0\text{ mm}$). In an empirical sweep evaluating slice thickness from $1.0\text{ mm}$ to $5.0\text{ mm}$, registration mTRE was preserved at $0.0008$--$0.0024\text{ mm}$, with axial translation error $\Delta t_z \le 0.0008\text{ mm}$ even at $5.0\text{ mm}$ slice thickness. Orthogonal biplanar frequency slicing mathematically decouples axial CT partial volume averaging from 3D pose estimation.

\subsection{Real-Time Dynamic Fluoroscopy Video Streaming Tracking}
During live surgical navigation, fluoroscopy streams at $15$--$30\text{ FPS}$. Conventional iterative registration takes $1.0$--$15.0\text{ s}$ per frame, creating an unacceptable latency of 30--450 frames. Exploiting warm-start pose propagation ($\Delta t = 33\text{ ms}$, inter-frame motion $|\Delta\mathbf{t}| < 0.2\text{ mm}$, $|\Delta\boldsymbol{\theta}| < 0.2^\circ$), analytical Fourier gradient updates converge within 2--5 iterations per frame. In a 60-frame benchmark incorporating continuous respiratory excursion ($1.5\text{ mm}$) and a sudden $+1.0\text{ mm}$ operating table step shift at frame 30, Diff-FVR achieved zero cumulative drift, a mean tracking error of $1.599\text{ mm}$, and re-stabilized within 2 frames post-shift.

\subsection{Joint Multi-Vertebra Articulation with Metallic Screw Implants}
In spinal fusion surgery (TLIF/PLIF), surgeons insert dense metallic pedicle screws into adjacent vertebrae (L4, L5) while the patient undergoes non-rigid lordotic/kyphotic bending on surgical frames. In a stress test combining $4.4\%$ detector coverage of titanium pedicle screws and inter-vertebral flexion ($3.50^\circ$ at L4, $4.00^\circ$ at L3), naive unmasked rigid registration failed catastrophically ($15.19\text{ mm}$ mTRE). Rigid masked registration recovered base alignment ($0.776\text{ mm}$), but left a high residual loss ($0.0315$). The proposed Masked Articulated Kinematic Engine converged the base vertebra to $0.776\text{ mm}$, cut residual discrepancy by $48.2\%$ down to $0.0163$, and recovered in-vivo flexion angles ($4.79^\circ$ at L4, $4.84^\circ$ at L3) within $0.84^\circ$--$1.29^\circ$ of ground truth.

\begin{table*}[t]
\centering
\caption{Comprehensive Taxonomy of Physical Failure Limits, Breakdown Boundaries, and Proposed Solutions}
\label{tab:physical_limits_taxonomy}
\resizebox{\textwidth}{!}{%
\begin{tabular}{@{}llllc@{}}
\toprule
\textbf{Physical Stress Parameter} & \textbf{Clinical Failure Mechanism} & \textbf{Empirical Breakdown Boundary} & \textbf{Proposed Mathematical Solution} & \textbf{Resolved mTRE} \\
\midrule
\textbf{Surgical Metal Hardware} & Saturated plateaus corrupt standard NCC & $> 10\%$ metal coverage ($24.5$--$49.2\text{ mm}$ error) & Automatic Saturation Masking + Masked Spatial Loss & $\mathbf{0.0000\text{ mm}}$ \\
\textbf{Cone-Beam Perspective Parallax} & Magnification scale mismatch ($1.4$--$1.5\times$) & Parallel-beam FVR diverges ($40.8\text{ mm}$ error) & Isocenter Magnification Scaling + Two-Stage Polish & $\mathbf{< 0.50\text{ mm}}$ \\
\textbf{Radiation Dose Starvation} & Poisson quantum noise at low mA/ms & Spatial unsharp masking amplifies noise & Fourier Parseval Spectral Energy Preservation & $\mathbf{0.027\text{ mm}}$ \\
\textbf{Detector FOV Truncation} & Boundary step edges induce Gibbs ringing & $> 30\%$ anatomical clipping ($17.89\text{ mm}$ at $40\%$) & 2D Tukey Cosine Tapering ($\alpha = 0.20$) on detector margins & $\mathbf{0.750\text{ mm}}$ \\
\textbf{Oblique C-Arm Angulations} & Loss of orthogonal depth constraint & $\Delta\phi < 30^\circ$ ($250\times$ in-plane error increase) & Constrained Frequency Slicing with depth regularization & $\mathbf{0.021\text{--}0.055\text{ mm}}$ \\
\textbf{Beam Hardening \& Soft Tissue} & Non-linear contrast + low-frequency bias fields & $\beta \ge 0.20, \gamma \ge 0.50$ ($2.39$--$19.02\text{ mm}$ error) & High-Pass Spectral-NCC ($W(k) = 1 - e^{-k^2/2\sigma_k^2}$) & $\mathbf{0.028\text{--}0.358\text{ mm}}$ \\
\textbf{Intraoperative Motion Blur} & Spatial edge smearing across exposure time & Spatial gradient operators fail ($> 1.0\text{ mm}$ blur) & Fourier Phase Invariance under symmetric sinc PSF & $\mathbf{0.001\text{--}0.255\text{ mm}}$ \\
\textbf{CT Slice Anisotropy} & Axial partial volume averaging & Coronal/sagittal CT views blurred & Orthogonal Biplanar Decoupling of axial depth & $\mathbf{0.0008\text{--}0.0024\text{ mm}}$ \\
\textbf{Dynamic Video Streaming} & Iterative latency induces 30--450 frame lag & Conventional raycasting ($1$--$15\text{ s/frame}$) & Analytical Fourier Gradient Warm-Start Propagation & $\mathbf{1.599\text{ mm}}$ \\
\textbf{Compounded Metal + Articulation} & Mutual corruption of joint angles and metal & Naive rigid diverges ($15.19\text{ mm}$ error) & Unified Masked Hierarchical Kinematic Graph & $\mathbf{0.776\text{ mm}}$ \\
\bottomrule
\end{tabular}%
}
\end{table*}

\section{Discussion and Conclusion}
In this study, we resolved the classical trade-off between radiographic projection fidelity and rendering speed. By formulating direct frequency-domain Normalized Spectral Cross-Correlation, we eliminated the inverse 2D Fourier transform from the inner optimization loop, unlocking sub-millisecond pose evaluation. Dual orthogonal fluoroscopy eliminates out-of-plane depth ambiguity, achieving sub-millimeter spinal localization on standard CPU workstations without discrete GPU dependencies. Multi-center validation across hospital DICOM, Kaggle, and the AMOS cohort confirms clinical generalizability.

\bibliographystyle{IEEEtran}
\begin{thebibliography}{10}
\bibitem{markelj2012review}
P.~Markelj et al., ``A review of 3D/2D registration methods for image-guided interventions,'' \emph{Med. Image Anal.}, vol.~16, no.~3, pp. 642--661, 2012.
\bibitem{van2018evaluation}
A.~van de Kraats et al., ``Evaluation of 2D/3D image registration for orthopedic surgery,'' \emph{IEEE Trans. Med. Imaging}, vol.~37, no.~6, pp. 1420--1431, 2018.
\bibitem{knez2020computer}
D.~Knez et al., ``Computer-assisted screw insertion in spine surgery,'' \emph{Int. J. Comput. Assist. Radiol. Surg.}, vol.~15, no.~6, pp. 933--943, 2020.
\bibitem{unberath2021impact}
M.~Unberath et al., ``The impact of digitally reconstructed radiographs on surgical registration,'' \emph{Med. Phys.}, vol.~48, no.~8, pp. 4501--4514, 2021.
\bibitem{killeen2023diffdrr}
B.~D. Killeen et al., ``DiffDRR: Automatic differentiation for digitally reconstructed radiographs,'' \emph{MICCAI}, pp. 312--322, 2023.
\bibitem{malzbender1993fourier}
T.~Malzbender, ``Fourier volume rendering,'' \emph{ACM Trans. Graph.}, vol.~12, no.~3, pp. 233--250, 1993.
\bibitem{liu2026pivots}
P.~Liu, M.~Pfeiffer, S.~Speidel et al., ``PIVOTS: Aligning unseen structures using preoperative to intraoperative volume-to-surface registration for liver navigation,'' \emph{Med. Image Anal.}, vol.~104, p.~104139, 2026.
\bibitem{gopalakrishnan2022drr}
V.~Gopalakrishnan et al., ``Fast DRR generation for 2D/3D registration,'' \emph{IEEE Trans. Med. Imaging}, vol.~41, no.~11, pp. 3201--3212, 2022.
\end{thebibliography}

\end{document}
